\documentclass[12pt,a4paper]{article}
\usepackage[margin=25mm, headheight=15pt, headsep=10pt]{geometry}
\usepackage{fontspec}
\usepackage[table]{xcolor} % Note: table option is required for rowcolors
\usepackage{titlesec}
\usepackage{tocloft}
\usepackage{fancyhdr}
\usepackage{booktabs}
\usepackage{tcolorbox}
\tcbuselibrary{skins, breakable}
\usepackage[colorlinks=true, linkcolor=emerald, urlcolor=emerald, bookmarks=true]{hyperref}

\definecolor{emerald}{RGB}{0,102,51}
\definecolor{darkred}{RGB}{139,0,0}
\definecolor{lightgray}{RGB}{245,245,245}

\usepackage[nil]{babel}
\babelprovide[import, main]{bengali}
\babelprovide[import]{arabic}
\babelfont{rm}[Renderer=HarfBuzz,Scale=1.0]{Noto Serif Bengali}
\babelfont{sf}[Renderer=HarfBuzz]{Noto Sans Bengali}
\babelfont{tt}[Renderer=HarfBuzz]{Noto Sans Bengali}
\babelfont[arabic]{rm}[Renderer=HarfBuzz,Scale=1.1]{Noto Naskh Arabic}

% Section styling
\titleformat{\section}{\Large\bfseries\color{emerald}}{\thesection.}{0.5em}{}
\titleformat{\subsection}{\large\bfseries\color{emerald!85!black}}{\thesubsection.}{0.5em}{}
\titleformat{\subsubsection}{\normalsize\bfseries\color{emerald!70!black}}{\thesubsubsection.}{0.5em}{}
\titlespacing*{\section}{0pt}{18pt}{8pt}

% Fancy Header/Footer setup
\pagestyle{fancy}
\fancyhf{} % clear all header and footer fields
\fancyhead[L]{\color{emerald}\small\textbf{\nouppercase{\leftmark}}}
\fancyfoot[L]{\scriptsize\color{black!50}{\lat{AI}}-সংকলিত, আলেম-যাচাইকৃত নয়}
\fancyfoot[C]{\color{emerald!70!black}\thepage}
\renewcommand{\headrulewidth}{0.4pt}
\renewcommand{\headrule}{\hbox to\headwidth{\color{emerald}\leaders\hrule height \headrulewidth\hfill}}
\renewcommand{\footrulewidth}{0pt}

% Custom boxes
% 1. Ayat/Hadith Box
\newtcolorbox{ayatbox}[1][]{
    enhanced, breakable, 
    colback=emerald!5, 
    colframe=emerald, 
    boxrule=0.5pt, 
    arc=3pt, 
    left=10pt, right=10pt, top=10pt, bottom=10pt,
    #1
}

% 2. Quote/Translation Box
\newtcolorbox{quotebox}[1][]{
    enhanced, breakable,
    frame hidden,
    colback=lightgray,
    borderline west={3pt}{0pt}{emerald},
    left=15pt, right=10pt, top=10pt, bottom=10pt,
    sharp corners,
    #1
}

% 3. AI-generated notice box
\newtcolorbox{warnbox}[1][]{
    enhanced, breakable,
    colback=red!4,
    colframe=darkred,
    boxrule=0.8pt,
    arc=3pt,
    left=10pt, right=10pt, top=10pt, bottom=10pt,
    #1
}

% Commands
\newcommand{\arab}[1]{\foreignlanguage{arabic}{#1}}
\newcommand{\kib}[1]{\textcolor{emerald}{\textbf{#1}}}

% Latin (English) fallback: Noto Serif Bengali has NO Latin glyphs
\newfontfamily\latfont[Renderer=HarfBuzz]{Noto Serif}
\newcommand{\lat}[1]{{\latfont #1}}

\setlength{\parskip}{5pt}
\setlength{\parindent}{0pt}

% Fix for missing bullet in Noto Serif Bengali
\renewcommand{\labelitemi}{$\bullet$}



\begin{document}
\begin{center}{\Large\bfseries\color{emerald} \lat{02}-সানা-ইস্তিফতাহ}\end{center}
\vspace{1em}
\section*{০২ — সানা (ইস্তিফতাহ): \textit{\lat{Subhanaka} \lat{Allahumma}...} — নামাজের শুরুর দোয়া}

\begin{warnbox}
 \textbf{গুরুত্বপূর্ণ নোটিশ (\lat{Important} \lat{Notice}):} এই কন্টেন্টটি \lat{AI} (কৃত্রিম বুদ্ধিমত্তা) দিয়ে প্রস্তুত ও সংকলিত। \lat{AI} কঠোর সূত্র-মান নীতি মেনে শুধু নির্ভরযোগ্য উৎস থেকে তথ্য সংগ্রহ করেছে — কেবল সহীহ/হাসান হাদিস ও সহীহ সনদের আসার রাখা হয়েছে, দুর্বল সনদ বাদ দেওয়া হয়েছে। তবে এটি কোনো আলেম বা মুহাদ্দিস কর্তৃক যাচাইকৃত নয়।
\end{warnbox}

\medskip\hrule\medskip

\subsection*{১. সূত্র কার্ড}

\begin{table}[h]\centering\small
\begin{tabular}{ll}
বিষয় & বিবরণ \\
\hline
\textbf{নাম} & সানা (ইস্তিফতাহ) — তাকবীরের পরের দোয়া \\
\textbf{সূত্র} & সুনানে আবু দাউদ ৭৭৫, সুনানে তিরমিযী ২৪৩ (সহীহ — আলবানি) \\
\textbf{মর্যাদা} & \textbf{সুন্নাতে মুআক্কাদাহ} — তাকবীরে তাহরিমার পর, ফাতিহার আগে \\
\hline
\end{tabular}
\end{table}

\medskip\hrule\medskip

\subsection*{২. মূল আরবি}

\begin{center}\arab{سُبْحَانَكَ اللَّهُمَّ وَبِحَمْدِكَ وَتَبَارَكَ اسْمُكَ وَتَعَالَىٰ جَدُّكَ وَلَا إِلَٰهَ غَيْرُكَ}\end{center}

\subsection*{৩. বাংলা উচ্চারণ}

\textbf{সুবহানাকাল্লাহুম্মা ওয়া বিহামদিকা ওয়া তাবারাকাস মুকা ওয়া তা'আলা জাদ্দুকা ওয়ালা ইলাহা গাইরুকা}

\medskip\hrule\medskip

\subsection*{৪. শব্দ-শব্দ অর্থ}

\begin{table}[h]\centering\small
\begin{tabular}{lll}
আরবি & বাংলা & \lat{English} \\
\hline
\textbf{\arab{سُبْحَانَكَ}} & আমি তোমাকে পবিত্র ঘোষণা করি & \textit{\lat{Glorified} \lat{be} \lat{You}} \\
\textbf{\arab{اللَّهُمَّ}} & হে আল্লাহ & \textit{\lat{O} \lat{Allah}} \\
\textbf{\arab{وَبِحَمْدِكَ}} & ও তোমার প্রশংসায় & \textit{\lat{and} \lat{with} \lat{Your} \lat{praise}} \\
\textbf{\arab{وَتَبَارَكَ}} & ও তোমার নাম বরকতময় & \textit{\lat{and} \lat{blessed} \lat{is}} \\
\textbf{\arab{اسْمُكَ}} & তোমার নাম & \textit{\lat{Your} \lat{Name}} \\
\textbf{\arab{وَتَعَالَىٰ}} & ও মহান & \textit{\lat{and} \lat{exalted} \lat{is}} \\
\textbf{\arab{جَدُّكَ}} & তোমার মর্যাদা/গৌরব & \textit{\lat{Your} \lat{Majesty}} \\
\textbf{\arab{وَلَا} \arab{إِلَٰهَ}} & ও কোনো ইলাহ নেই & \textit{\lat{and} \lat{there} \lat{is} \lat{no} \lat{god}} \\
\textbf{\arab{غَيْرُكَ}} & তোমার বাইরে & \textit{\lat{besides} \lat{You}} \\
\hline
\end{tabular}
\end{table}

\medskip\hrule\medskip

\subsection*{৫. পূর্ণ বাংলা অনুবাদ}

\textbf{\lat{“}হে আল্লাহ! আমি তোমাকে পবিত্র ঘোষণা করি এবং তোমার প্রশংসা করি। তোমার নাম বরকতময়, তোমার মর্যাদা মহান — তোমার বাইরে কোনো ইলাহ নেই।\lat{”}}

\medskip\hrule\medskip

\subsection*{৬. সংক্ষিপ্ত ব্যাখ্যা (\lat{Khushu-focused})}

\subsubsection*{\textbf{"সুবহানাকা" — পবিত্রতা ঘোষণা}}
"\lat{Subhanaka}" আল্লাহকে \lat{pure} (বিশুদ্ধ) ঘোষণা করা — সকল \lat{imperfection} (অপূর্ণতা) থেকে \lat{free} (মুক্ত), সকল \lat{perfection} (পূর্ণতা)-র উৎস।

\subsubsection*{\textbf{"বিহামদিকা" — প্রশংসায়}}
"\lat{Bihamdika}" — তোমার প্রশংসায়। শুধু \lat{Subhanaka} যথেষ্ট নয় — প্রশংসা (\lat{hamd}) যোগ করা হয়েছে। এটি \lat{gratitude} (কৃতজ্ঞতা) ও \lat{acknowledgment} (স্বীকৃতি)।

\subsubsection*{\textbf{"তাবারাকাস মুকা" — বরকত}}
"\lat{Tabarakasmuka}" — তোমার নাম বরকতময়। \lat{Blessing} (বরকত) মানে: যা থেকে নিরন্তর কল্যাণ বর্ষিত হয় — \lat{sustenance} (জীবিকা), \lat{guidance} (হেদায়াত), \lat{mercy} (রহমাত)।

\subsubsection*{\textbf{"তা'আলা জাদ্দুকা" — গৌরব}}
"\lat{Ta'ala} \lat{Jadduka}" — তোমার মর্যাদা মহান। \lat{Majesty} (মর্যাদা) ও \lat{glory} (গৌরব)।

\subsubsection*{\textbf{"ওয়ালা ইলাহা গাইরুকা" — তাওহীদ}}
"\lat{La} \lat{ilaha} \lat{ghairuka}" — তোমার বাইরে কোনো ইলাহ নেই। \lat{Tawhid} (তাওহীদ)-এর সার।

\medskip\hrule\medskip

\subsection*{৭. খুশুর জন্য মূল পয়েন্টস}

\begin{table}[h]\centering\small
\begin{tabular}{ll}
পয়েন্ট & কীভাবে প্রয়োগ করবেন \\
\hline
\textbf{গভীর অর্থ} & প্রতিটি শব্দের অর্থ মনে করুন — "আমি আল্লাহকে পবিত্র ঘোষণা করছি" \\
\textbf{হাতের অবস্থান} & তাকবীরের পর হাত বাঁধুন — দায়েন বামের ওপর (শাফি'ি) / নাভির নিচে (হানাফী) \\
\textbf{চোখের অবস্থান} & কিবলার দিকে স্থির — ভুমির দিকে তাকিয়ে \\
\textbf{গলার স্বর} & নিজে শোনার মতো ধীরে বলুন — অন্যকে নয় \\
\hline
\end{tabular}
\end{table}

\medskip\hrule\medskip

\subsection*{৮. সংশ্লিষ্ট হাদিস}

\begin{table}[h]\centering\small
\begin{tabular}{lll}
বিষয় & হাদিস & গ্রন্থ \\
\hline
সানা পড়ার বিষয় & \textbf{\lat{“}তাকবীরের পর তোমাদের প্রত্যেকে 'সুবহানাকাল্লাহুম্মা' পড়ুক।\lat{”}} & সুনানে আবু দাউদ ৭৭৫ \\
সানা ও তাকবীরের সময় & \textbf{\lat{“}তোমরা নামাজে 'আল্লাহু আকবার' বলো, এরপর 'সুবহানাকাল্লাহুম্মা' পড়ো।\lat{”}} & সুনানে তিরমিযী ২৪৩ \\
\hline
\end{tabular}
\end{table}

\medskip\hrule\medskip

\textit{শেষ আপডেট: আগস্ট ২০২৬}


\end{document}
